\documentclass[11pt,a4paper]{article}

%% ── Packages ──────────────────────────────────────────────────────
\usepackage[utf8]{inputenc}
\usepackage[T1]{fontenc}
\usepackage{lmodern}
\usepackage{amsmath, amsthm, amssymb}
\usepackage{hyperref}
\usepackage{xcolor}
\usepackage{cleveref}
\usepackage{microtype}

%% ── Global typesetting tolerances ────────────────────────────────
\tolerance=1000
\emergencystretch=1.5em
\hbadness=10000

%% ── Lean code-listing support ────────────────────────────────────
\input{lean-setup}

%% ── Title block ──────────────────────────────────────────────────
\title{Unified Algebraic-Lattice Bipartition Framework:\\
       A Formal Approach}
\author{Frederick de Ruiter}
\date{\today}

\begin{document}
\newcommand{\TelemetryPhaseTwoTime}{5.00}
\newcommand{\TelemetryPhaseTwoBranches}{1,000}
\newcommand{\TelemetryPruned}{200}
\newcommand{\TelemetryMaxLog}{43}
\newcommand{\TelemetryMinLog}{37}
\newcommand{\TelemetryCertHash}{0ce4fa27fa33}
\newcommand{\TelemetryPhaseOnePruned}{0}
\newcommand{\TelemetryTotalTime}{0 hours, 0 minutes, 5 seconds}
\newcommand{\TelemetryPhaseOneTime}{0 hours, 0 minutes, 0 seconds}
\newcommand{\TelemetryNodesPerSec}{200}
\newcommand{\TelemetryAbundancePct}{100.0}
\newcommand{\TelemetryRaycastPct}{0.0}
\newcommand{\TelemetryEngineVersion}{1.0.0}
\newcommand{\TelemetryCommitHash}{3528dcbbbcebec32a81a5be431299cfffc01226d}
\newcommand{\TelemetryBoundsEnforced}{True}
\newcommand{\TelemetryHagisBaselineMinPrimeFactors}{7}
\newcommand{\TelemetryPrasadSunithaBound}{15}
\newcommand{\HashUALBFEngineRuleASafe}{ecc9b4ff5e44c09c8d5dca07b2cab6259bde792d25b5c269aad804bf417a3894}
\newcommand{\HashUALBFEngineRuleASafeStatus}{proven}
\newcommand{\HashUALBFEngineRuleBSafe}{ecc9b4ff5e44c09c8d5dca07b2cab6259bde792d25b5c269aad804bf417a3894}
\newcommand{\HashUALBFEngineRuleBSafeStatus}{proven}
\newcommand{\HashUALBFEnginePrasadSunithaComponentBound}{ecc9b4ff5e44c09c8d5dca07b2cab6259bde792d25b5c269aad804bf417a3894}
\newcommand{\HashUALBFEnginePrasadSunithaComponentBoundStatus}{proven}
\newcommand{\HashUALBFEngineCyclotomicGraphForcedInclusion}{98af5f79cfd707745092ef69f60071ca6dc923a1b84e4e448e0701ddff2791c9}
\newcommand{\HashUALBFEngineCyclotomicGraphForcedInclusionStatus}{proven}
\newcommand{\HashUALBFEngineCyclotomicGraphTransitiveForcedInclusion}{98af5f79cfd707745092ef69f60071ca6dc923a1b84e4e448e0701ddff2791c9}
\newcommand{\HashUALBFEngineCyclotomicGraphTransitiveForcedInclusionStatus}{proven}
\newcommand{\HashUALBFEngineCyclotomicGraphTransitiveReachabilitySoundness}{98af5f79cfd707745092ef69f60071ca6dc923a1b84e4e448e0701ddff2791c9}
\newcommand{\HashUALBFEngineCyclotomicGraphTransitiveReachabilitySoundnessStatus}{proven}
\newcommand{\HashUALBFEngineSieveSoundnessRustSieveSoundness}{02184c7e367704a945ce5c7f36d76f26104e92b0619dabbfcdb0eeee8381b1a5}
\newcommand{\HashUALBFEngineSieveSoundnessRustSieveSoundnessStatus}{proven}
\newcommand{\HashUALBFEngineBipartitionPrefixSigmaCoprime}{aab92a155825b8cde3047488d215b64207fa8a855e7799024ae2b7a77d02ff76}
\newcommand{\HashUALBFEngineBipartitionPrefixSigmaCoprimeStatus}{proven}
\newcommand{\HashUALBFEngineBipartitionAmbsSuffixTarget}{aab92a155825b8cde3047488d215b64207fa8a855e7799024ae2b7a77d02ff76}
\newcommand{\HashUALBFEngineBipartitionAmbsSuffixTargetStatus}{proven}
\newcommand{\HashUALBFEngineBipartitionNoSolutionNoQpn}{aab92a155825b8cde3047488d215b64207fa8a855e7799024ae2b7a77d02ff76}
\newcommand{\HashUALBFEngineBipartitionNoSolutionNoQpnStatus}{proven}
\newcommand{\HashUALBFQPNAbundancyBoundQpnAbundancyTarget}{a13219893d7bf6ba12cef6062a2351959fcfe522446ac1cba0aa462bb4ed3fe3}
\newcommand{\HashUALBFQPNAbundancyBoundQpnAbundancyTargetStatus}{proven}
\newcommand{\HashUALBFQPNAbundancyBoundQpnTotientBound}{a13219893d7bf6ba12cef6062a2351959fcfe522446ac1cba0aa462bb4ed3fe3}
\newcommand{\HashUALBFQPNAbundancyBoundQpnTotientBoundStatus}{proven}
\newcommand{\HashUALBFQPNAbundancyBoundAbundancyStarvation}{a13219893d7bf6ba12cef6062a2351959fcfe522446ac1cba0aa462bb4ed3fe3}
\newcommand{\HashUALBFQPNAbundancyBoundAbundancyStarvationStatus}{proven}
\newcommand{\HashUALBFQPNObstructionLegendreCattaneoObstruction}{3dbfc95f27c80270af2610fc15c4df57a341f814453e3dd307cb09fb4e1cfb5d}
\newcommand{\HashUALBFQPNObstructionLegendreCattaneoObstructionStatus}{proven}
\newcommand{\HashUALBFQPNBasicPropertiesQpnIsOddSquare}{8992defd0cff60fce435d433d20749c86bdc42ca822648993c86c1487a29abb6}
\newcommand{\HashUALBFQPNBasicPropertiesQpnIsOddSquareStatus}{proven}
\newcommand{\HashUALBFQPNPrasadSunithaQpnCoprimeOneFiveOmegaBound}{84c169b594060172704bc94ffcdf95e99dce6ec0221279e2d4e632edcb0939c8}
\newcommand{\HashUALBFQPNPrasadSunithaQpnCoprimeOneFiveOmegaBoundStatus}{proven}
\newcommand{\HashUALBFQPNPrasadSunithaQpnDivFiveCoprimeThreeOmegaBound}{84c169b594060172704bc94ffcdf95e99dce6ec0221279e2d4e632edcb0939c8}
\newcommand{\HashUALBFQPNPrasadSunithaQpnDivFiveCoprimeThreeOmegaBoundStatus}{proven}
\newcommand{\HashUALBFEngineObstructionQpnSigmaModThree}{7b1260c280195bb3b99caa4267fdb6dfe13cc0e9b9e3f78e490a1f78c19307e7}
\newcommand{\HashUALBFEngineObstructionQpnSigmaModThreeStatus}{proven}
\newcommand{\HashUALBFEngineObstructionQpnSigmaModNine}{7b1260c280195bb3b99caa4267fdb6dfe13cc0e9b9e3f78e490a1f78c19307e7}
\newcommand{\HashUALBFEngineObstructionQpnSigmaModNineStatus}{proven}
\newcommand{\HashUALBFQPNTouchardQPNQpnSigmaModTwoFour}{9cc566426952947f41130b267e82abef3dd8d539ca1a1e477c0c6a0239230a25}
\newcommand{\HashUALBFQPNTouchardQPNQpnSigmaModTwoFourStatus}{proven}
\newcommand{\HashUALBFEngineTouchardBridgeTouchardBridge}{3ff39a6f0b2d24caf1e9da8832f14d1abc141f7b33ad14443e11b299bca5ec20}
\newcommand{\HashUALBFEngineTouchardBridgeTouchardBridgeStatus}{proven}
\newcommand{\HashUALBFFFIFromUFiveOneTwoToUFiveOneTwo}{a0bd1410394097b7ac0ed1d390274238959efde45008de146ea29356b614187d}
\newcommand{\HashUALBFFFIFromUFiveOneTwoToUFiveOneTwoStatus}{proven}
\newcommand{\HashUALBFFFIToUFiveOneTwoFromUFiveOneTwo}{a0bd1410394097b7ac0ed1d390274238959efde45008de146ea29356b614187d}
\newcommand{\HashUALBFFFIToUFiveOneTwoFromUFiveOneTwoStatus}{proven}
\newcommand{\HashUALBFFFIFromUFiveOneTwoEqFromUFiveOneTwoFast}{0d17fd36a53e645c80ef6c67c91a034110e2a2d6317a9fe44c6abe4b2fc07ba3}
\newcommand{\HashUALBFFFIFromUFiveOneTwoEqFromUFiveOneTwoFastStatus}{proven}
\newcommand{\HashUALBFFFIModInverseSpec}{a0bd1410394097b7ac0ed1d390274238959efde45008de146ea29356b614187d}
\newcommand{\HashUALBFFFIModInverseSpecStatus}{proven}
\newcommand{\HashUALBFFFIUalbfModInverseOkLimbsEq}{a0bd1410394097b7ac0ed1d390274238959efde45008de146ea29356b614187d}
\newcommand{\HashUALBFFFIUalbfModInverseOkLimbsEqStatus}{proven}
\newcommand{\HashUALBFFFIUalbfModInverseLimbEq}{a0bd1410394097b7ac0ed1d390274238959efde45008de146ea29356b614187d}
\newcommand{\HashUALBFFFIUalbfModInverseLimbEqStatus}{proven}
\newcommand{\HashUALBFFFIUalbfCheckCrtOneOneFiveFiveLimbsEq}{a0bd1410394097b7ac0ed1d390274238959efde45008de146ea29356b614187d}
\newcommand{\HashUALBFFFIUalbfCheckCrtOneOneFiveFiveLimbsEqStatus}{proven}
\newcommand{\HashUALBFFFIUFiveOneTwoWZeroMk}{0d17fd36a53e645c80ef6c67c91a034110e2a2d6317a9fe44c6abe4b2fc07ba3}
\newcommand{\HashUALBFFFIUFiveOneTwoWZeroMkStatus}{proven}
\newcommand{\HashUALBFFFIUFiveOneTwoWOneMk}{0d17fd36a53e645c80ef6c67c91a034110e2a2d6317a9fe44c6abe4b2fc07ba3}
\newcommand{\HashUALBFFFIUFiveOneTwoWOneMkStatus}{proven}
\newcommand{\HashUALBFFFIUFiveOneTwoWTwoMk}{0d17fd36a53e645c80ef6c67c91a034110e2a2d6317a9fe44c6abe4b2fc07ba3}
\newcommand{\HashUALBFFFIUFiveOneTwoWTwoMkStatus}{proven}
\newcommand{\HashUALBFFFIUFiveOneTwoWThreeMk}{0d17fd36a53e645c80ef6c67c91a034110e2a2d6317a9fe44c6abe4b2fc07ba3}
\newcommand{\HashUALBFFFIUFiveOneTwoWThreeMkStatus}{proven}
\newcommand{\HashUALBFFFIUFiveOneTwoWFourMk}{0d17fd36a53e645c80ef6c67c91a034110e2a2d6317a9fe44c6abe4b2fc07ba3}
\newcommand{\HashUALBFFFIUFiveOneTwoWFourMkStatus}{proven}
\newcommand{\HashUALBFFFIUFiveOneTwoWFiveMk}{0d17fd36a53e645c80ef6c67c91a034110e2a2d6317a9fe44c6abe4b2fc07ba3}
\newcommand{\HashUALBFFFIUFiveOneTwoWFiveMkStatus}{proven}
\newcommand{\HashUALBFFFIUFiveOneTwoWSixMk}{0d17fd36a53e645c80ef6c67c91a034110e2a2d6317a9fe44c6abe4b2fc07ba3}
\newcommand{\HashUALBFFFIUFiveOneTwoWSixMkStatus}{proven}
\newcommand{\HashUALBFFFIUFiveOneTwoWSevenMk}{0d17fd36a53e645c80ef6c67c91a034110e2a2d6317a9fe44c6abe4b2fc07ba3}
\newcommand{\HashUALBFFFIUFiveOneTwoWSevenMkStatus}{proven}
\newcommand{\HashUALBFPureABCConjectureDeriveConjecturalCeiling}{e22eaa1a38904f80693db2f12ceddc0e3c5b45a413341ae500a22c092ae69415}
\newcommand{\HashUALBFPureABCConjectureDeriveConjecturalCeilingStatus}{proven}
\newcommand{\HashUALBFPureABCConjectureConjecturalCeilingSizeExclusion}{e22eaa1a38904f80693db2f12ceddc0e3c5b45a413341ae500a22c092ae69415}
\newcommand{\HashUALBFPureABCConjectureConjecturalCeilingSizeExclusionStatus}{proven}
\newcommand{\HashUALBFEngineModOneOneFiveFiveBridgeModEqOfModEqOfDvd}{951a56c679bde4e396f206ad099f2998e808d8491de7a04327c3631a6d300f9d}
\newcommand{\HashUALBFEngineModOneOneFiveFiveBridgeModEqOfModEqOfDvdStatus}{proven}
\newcommand{\HashUALBFEngineModOneOneFiveFiveBridgeModOneOneFiveFiveToModThree}{951a56c679bde4e396f206ad099f2998e808d8491de7a04327c3631a6d300f9d}
\newcommand{\HashUALBFEngineModOneOneFiveFiveBridgeModOneOneFiveFiveToModThreeStatus}{proven}
\newcommand{\HashUALBFEngineModOneOneFiveFiveBridgeModOneOneFiveFiveToModFive}{951a56c679bde4e396f206ad099f2998e808d8491de7a04327c3631a6d300f9d}
\newcommand{\HashUALBFEngineModOneOneFiveFiveBridgeModOneOneFiveFiveToModFiveStatus}{proven}
\newcommand{\HashUALBFEngineModOneOneFiveFiveBridgeModOneOneFiveFiveToModSeven}{951a56c679bde4e396f206ad099f2998e808d8491de7a04327c3631a6d300f9d}
\newcommand{\HashUALBFEngineModOneOneFiveFiveBridgeModOneOneFiveFiveToModSevenStatus}{proven}
\newcommand{\HashUALBFEngineModOneOneFiveFiveBridgeModOneOneFiveFiveToModOneOne}{951a56c679bde4e396f206ad099f2998e808d8491de7a04327c3631a6d300f9d}
\newcommand{\HashUALBFEngineModOneOneFiveFiveBridgeModOneOneFiveFiveToModOneOneStatus}{proven}
\newcommand{\HashUALBFEngineModOneOneFiveFiveBridgeModOneOneFiveFiveSoundness}{951a56c679bde4e396f206ad099f2998e808d8491de7a04327c3631a6d300f9d}
\newcommand{\HashUALBFEngineModOneOneFiveFiveBridgeModOneOneFiveFiveSoundnessStatus}{proven}
\newcommand{\HashUALBFEngineModOneOneFiveFiveBridgeUalbfCheckCrtOneOneFiveFiveSound}{951a56c679bde4e396f206ad099f2998e808d8491de7a04327c3631a6d300f9d}
\newcommand{\HashUALBFEngineModOneOneFiveFiveBridgeUalbfCheckCrtOneOneFiveFiveSoundStatus}{proven}
\newcommand{\HashUALBFFixedSixFourScaleBoundCeilConservative}{c01161d7f14eb9d953a6908cb77497a53c2ae55d54f48a7271545bfd4a497ece}
\newcommand{\HashUALBFFixedSixFourScaleBoundCeilConservativeStatus}{proven}
\newcommand{\HashUALBFEngineSieveSoundnessRustSieveSoundnessModFive}{02184c7e367704a945ce5c7f36d76f26104e92b0619dabbfcdb0eeee8381b1a5}
\newcommand{\HashUALBFEngineSieveSoundnessRustSieveSoundnessModFiveStatus}{proven}
\newcommand{\HashUALBFEngineModFiveBridgeUalbfCheckModFiveSoundnessFfi}{eb01c656fa0f36dd4f5e713f683e58a7080201ad48acf389a226a42e03ec664f}
\newcommand{\HashUALBFEngineModFiveBridgeUalbfCheckModFiveSoundnessFfiStatus}{proven}
\newcommand{\HashUALBFEngineSieveSoundnessModularSieve}{02184c7e367704a945ce5c7f36d76f26104e92b0619dabbfcdb0eeee8381b1a5}
\newcommand{\HashUALBFEngineSieveSoundnessModularSieveStatus}{proven}
\newcommand{\HashUALBFQPNAbundancyBoundLeanAbundancyStarvationTheorem}{a13219893d7bf6ba12cef6062a2351959fcfe522446ac1cba0aa462bb4ed3fe3}
\newcommand{\HashUALBFQPNAbundancyBoundLeanAbundancyStarvationTheoremStatus}{proven}
\newcommand{\HashUALBFQPNAbundancyBoundVerifyStarvationPruning}{a13219893d7bf6ba12cef6062a2351959fcfe522446ac1cba0aa462bb4ed3fe3}
\newcommand{\HashUALBFQPNAbundancyBoundVerifyStarvationPruningStatus}{proven}
\newcommand{\HashUALBFQPNAbundancyBoundIsStarvedBound}{a13219893d7bf6ba12cef6062a2351959fcfe522446ac1cba0aa462bb4ed3fe3}
\newcommand{\HashUALBFQPNAbundancyBoundIsStarvedBoundStatus}{proven}
\newcommand{\HashUALBFEngineCyclotomicGraphIsCdgForcedPruned}{98af5f79cfd707745092ef69f60071ca6dc923a1b84e4e448e0701ddff2791c9}
\newcommand{\HashUALBFEngineCyclotomicGraphIsCdgForcedPrunedStatus}{proven}
\newcommand{\HashUALBFEngineBipartitionCoprimeMultiplicativeNonlinear}{aab92a155825b8cde3047488d215b64207fa8a855e7799024ae2b7a77d02ff76}
\newcommand{\HashUALBFEngineBipartitionCoprimeMultiplicativeNonlinearStatus}{proven}
\newcommand{\HashUALBFEngineBipartitionCoprimeMultiplicative}{aab92a155825b8cde3047488d215b64207fa8a855e7799024ae2b7a77d02ff76}
\newcommand{\HashUALBFEngineBipartitionCoprimeMultiplicativeStatus}{proven}
\newcommand{\HashUALBFEngineBipartitionDisjointByConstruction}{aab92a155825b8cde3047488d215b64207fa8a855e7799024ae2b7a77d02ff76}
\newcommand{\HashUALBFEngineBipartitionDisjointByConstructionStatus}{proven}
\newcommand{\HashUALBFQPNPrasadSunithaVerifyPrasadSunitha}{84c169b594060172704bc94ffcdf95e99dce6ec0221279e2d4e632edcb0939c8}
\newcommand{\HashUALBFQPNPrasadSunithaVerifyPrasadSunithaStatus}{proven}
\newcommand{\HashUALBFQPNTouchardQPNScreenModEight}{9cc566426952947f41130b267e82abef3dd8d539ca1a1e477c0c6a0239230a25}
\newcommand{\HashUALBFQPNTouchardQPNScreenModEightStatus}{proven}
\newcommand{\HashUALBFQPNObstructionIsValidModEight}{3dbfc95f27c80270af2610fc15c4df57a341f814453e3dd307cb09fb4e1cfb5d}
\newcommand{\HashUALBFQPNObstructionIsValidModEightStatus}{proven}
\newcommand{\HashUALBFEngineSieveSoundnessPassesRaycastSieveSpec}{02184c7e367704a945ce5c7f36d76f26104e92b0619dabbfcdb0eeee8381b1a5}
\newcommand{\HashUALBFEngineSieveSoundnessPassesRaycastSieveSpecStatus}{proven}
\newcommand{\HashUALBFEngineSieveSoundnessVerifiedPassesRaycastSieve}{02184c7e367704a945ce5c7f36d76f26104e92b0619dabbfcdb0eeee8381b1a5}
\newcommand{\HashUALBFEngineSieveSoundnessVerifiedPassesRaycastSieveStatus}{proven}
\newcommand{\HashUALBFEngineCyclotomicGraphZsigmondyPreconditionsSatisfied}{98af5f79cfd707745092ef69f60071ca6dc923a1b84e4e448e0701ddff2791c9}
\newcommand{\HashUALBFEngineCyclotomicGraphZsigmondyPreconditionsSatisfiedStatus}{proven}
\newcommand{\HashUALBFEngineCyclotomicGraphProofVerifyZsigmondyPreconditions}{98af5f79cfd707745092ef69f60071ca6dc923a1b84e4e448e0701ddff2791c9}
\newcommand{\HashUALBFEngineCyclotomicGraphProofVerifyZsigmondyPreconditionsStatus}{proven}
\newcommand{\HashUALBFPureArithmeticLemmaCompositeHasPrimeFactorLeSqrt}{e272303ea6a318ffe7560f6a8d62cddc3d8903680dadf7f5563c4849ac043630}
\newcommand{\HashUALBFPureArithmeticLemmaCompositeHasPrimeFactorLeSqrtStatus}{proven}
\newcommand{\HashUALBFPureArithmeticLemmaSmallestFactorIsPrime}{e272303ea6a318ffe7560f6a8d62cddc3d8903680dadf7f5563c4849ac043630}
\newcommand{\HashUALBFPureArithmeticLemmaSmallestFactorIsPrimeStatus}{proven}
\newcommand{\HashUALBFPureArithmeticLemmaModpowModDivisibility}{e272303ea6a318ffe7560f6a8d62cddc3d8903680dadf7f5563c4849ac043630}
\newcommand{\HashUALBFPureArithmeticLemmaModpowModDivisibilityStatus}{proven}
\newcommand{\HashUALBFPureArithmeticLemmaModpowAddMul}{e272303ea6a318ffe7560f6a8d62cddc3d8903680dadf7f5563c4849ac043630}
\newcommand{\HashUALBFPureArithmeticLemmaModpowAddMulStatus}{proven}
\newcommand{\HashUALBFPureArithmeticLemmaOrderExists}{e272303ea6a318ffe7560f6a8d62cddc3d8903680dadf7f5563c4849ac043630}
\newcommand{\HashUALBFPureArithmeticLemmaOrderExistsStatus}{proven}
\newcommand{\HashUALBFPureArithmeticLemmaOrderPrimeFactor}{e272303ea6a318ffe7560f6a8d62cddc3d8903680dadf7f5563c4849ac043630}
\newcommand{\HashUALBFPureArithmeticLemmaOrderPrimeFactorStatus}{proven}
\newcommand{\HashUALBFPureArithmeticLemmaDivisibilityBounds}{e272303ea6a318ffe7560f6a8d62cddc3d8903680dadf7f5563c4849ac043630}
\newcommand{\HashUALBFPureArithmeticLemmaDivisibilityBoundsStatus}{proven}
\newcommand{\HashUALBFPureArithmeticLemmaFermatLittleTheorem}{e272303ea6a318ffe7560f6a8d62cddc3d8903680dadf7f5563c4849ac043630}
\newcommand{\HashUALBFPureArithmeticLemmaFermatLittleTheoremStatus}{proven}
\newcommand{\HashUALBFPureArithmeticLemmaOrderLePMinusOne}{e272303ea6a318ffe7560f6a8d62cddc3d8903680dadf7f5563c4849ac043630}
\newcommand{\HashUALBFPureArithmeticLemmaOrderLePMinusOneStatus}{proven}
\newcommand{\HashUALBFPureArithmeticLemmaSquareComparisonContradiction}{e272303ea6a318ffe7560f6a8d62cddc3d8903680dadf7f5563c4849ac043630}
\newcommand{\HashUALBFPureArithmeticLemmaSquareComparisonContradictionStatus}{proven}
\newcommand{\HashUALBFPureArithmeticLemmaFSquaredGtNMinusOne}{e272303ea6a318ffe7560f6a8d62cddc3d8903680dadf7f5563c4849ac043630}
\newcommand{\HashUALBFPureArithmeticLemmaFSquaredGtNMinusOneStatus}{proven}
\newcommand{\HashUALBFPureArithmeticLemmaPocklingtonCertificate}{e272303ea6a318ffe7560f6a8d62cddc3d8903680dadf7f5563c4849ac043630}
\newcommand{\HashUALBFPureArithmeticLemmaPocklingtonCertificateStatus}{proven}
\newcommand{\HashUALBFPureArithmeticLemmaDivisibilityTransitive}{e272303ea6a318ffe7560f6a8d62cddc3d8903680dadf7f5563c4849ac043630}
\newcommand{\HashUALBFPureArithmeticLemmaDivisibilityTransitiveStatus}{proven}
\newcommand{\HashCheckCdgForcedKill}{fd5dceec55e366bafd28b764e97aa3b84b9d29653293695d4bb32059f1da07b5}
\newcommand{\HashCheckStarvationKill}{5b9cbf43d6d7495fa265f2b0b6cb6d738586441f56d4ae2ccbcc2838e26187eb}
\newcommand{\HashCyclotomicEvalSpec}{25cebf82f225d163bd96469454c1e3ee657da89962649268029c4739cf9824e5}
\newcommand{\HashCyclotomicEvalSpecAux}{5312ec6bd2fffba34f9d3f45093729673e1d105325afbdc9ea8d15b6f6b19ff2}
\newcommand{\HashDisjointFactorSets}{67f57ef0b67f0ea8774630767795705c4cd897014ff82f26178d6358902ddae2}
\newcommand{\HashIsCdgForcedPruned}{4d58fa97e098febed1cc1e944e35fc603c9f404cb7955ba9c76453c1ddbf5ec9}
\newcommand{\HashIsOrderModP}{52f041664e6cc1e8660cb4c44c14adc659a23a5675f45630ea705c1e96b1cddc}
\newcommand{\HashIsPrime}{3f1f4635405e8940cfb5b0df21aed8dc88b742db0826f9ca5d043bb0b0c00998}
\newcommand{\HashIsStarved}{13a88aad6528b0a9512c1fc20f045fe62b60f5fa7c4eaf6e3dabc52cf28279dc}
\newcommand{\HashIsValidModEight}{757e6e04f11d3e5456a8e7204eba8e3ebb2c1c0b7096542751dd63f5421d7a9d}
\newcommand{\HashLeanAbundancyStarvationTheorem}{3f4300e9de5e8e30ef5786aef95f510ca2fdd189fc24e155441bfa357c7ea366}
\newcommand{\HashLeanConjecturalActive}{054c951db92a7753e908c5cd07bcb6e0786b190883b40aeba83f6a9b2596677e}
\newcommand{\HashLeanConjecturalMaxLogOneZeroCeiling}{29b066ebf0cc2f9b689209b69800255bb08dd334cc397331cfb81d6e87534a49}
\newcommand{\HashLeanCrtModulusProduct}{6e5832a4ed56ed180f0fd47c09e3f6540386901b93bd40f71eae31d4602d0dd3}
\newcommand{\HashLeanDivFiveCoprimeThreeBound}{58c00b2c255de978e040a24c6ad10d6ba6f60e69aa1be2cac8e45bff593752e8}
\newcommand{\HashLeanDivFiveCoprimeThreeCombined}{fa9dbeb60e5695908a443f37ff747ebecc5e1ed70eff1282b571d69a88d34a11}
\newcommand{\HashLeanDivFiveCoprimeThreeOffset}{f97bddc5b282c6826aacf1c69576f0fb851ea8ce35ea9ff9c8efcebe7a9fd074}
\newcommand{\HashLeanHagisOneNineEightTwoCombined}{d799a3d05070ab9ea82be983f68620938998316dedece75959a8265099e7b7a6}
\newcommand{\HashLeanHagisOneNineEightTwoMinPrimeFactors}{20c7b1fa4b03ee962857f42db73fb2658b765ffe2935810f755355d3ddcf5d12}
\newcommand{\HashLeanHagisOneNineEightTwoOffset}{e9da6d1340b0e16439ceec6f79165acfb6c2ebb274cef23623e82f3646e67bd6}
\newcommand{\HashLeanMaxExponent}{90f8882992d3c1df4148d039141f0fd1d77cde9c7a46424688a4e0b35e24d4d9}
\newcommand{\HashLeanMillerRabinTwoZeroBaseSufficiency}{1ac089edc418d75c955d64c01db144353552261e617c5e9959ed72ce09d40ff0}
\newcommand{\HashLeanOverflowThresholdDen}{16bbefdf025ac1858c60781b7d42fe70b4699681d11e08edc9668b9a62576f29}
\newcommand{\HashLeanOverflowThresholdNum}{21b1ceb0bc445c1ee792e0bbea96786306768805cbbef3bc204e4d0f05398feb}
\newcommand{\HashLeanPollardRhoBatchSize}{007aa18ae65f2100b6355b3accc3f40aa4711492a40e76fd76ba4e6c0905358a}
\newcommand{\HashLeanPollardRhoIterationLimit}{eac577d8d9b117ca1ae6472fd279d32fe518748eeb07fbb5056b1de462df5c90}
\newcommand{\HashLeanPrasadSunithaBound}{bd65e545edb5520cae3fb57071af5c9811da73d4d61f6c3dac40766e66bfdb3d}
\newcommand{\HashLeanPrasadSunithaCombined}{3824e4ca05d6d2b9619b29936a63c956bc32ad7c4895d999f670c936151b01e6}
\newcommand{\HashLeanPrasadSunithaOffset}{b2724d4e6a8188f28b01abea55635334c2e2f633fd94abdcd448c1e5d6e525de}
\newcommand{\HashLeanPrefixStopThreshold}{5ccd4cbf319de1022eb785e24388033bb2a56d87f5fed2e8a567811bc525aac7}
\newcommand{\HashLeanPrimeSplitThreshold}{1ed6edcd1aa2e917abc95af8396545b9ae1cacfe0e87ce4742488d5593af4f1d}
\newcommand{\HashLeanQpnTotientBoundDen}{4a1455d4e899a204fe791f8e2afea8ad1deedfa393bfae7f4ff9f864fe2712d9}
\newcommand{\HashLeanQpnTotientBoundNum}{9e215c746fefba6057f18a1f9047b3391b35a17fe6fe321b51663ea498097685}
\newcommand{\HashLeanRaycastChunkSize}{28d2c3963d17ca81804e9dbd05cb824e022f811da8a7ee93b05e814fa6a321ce}
\newcommand{\HashLeanRaycastGpuThreshold}{3817ab7c3e029a5dcb458d278e651908c23ab2d3c67f5f66afb77c187eb8059b}
\newcommand{\HashLeanSieveLimit}{d8e81b91dd484458558f59b3404ef888ec701a2ddb92210d9757892bd1af6bd4}
\newcommand{\HashLeanTargetMaxLogOneZero}{938dc08d46d78e63e0b91559a735070ddf157e3ff4c944c92afc8d48694fa113}
\newcommand{\HashLeanTargetMinLogOneZero}{4d471ee6e66fe3c30f3313d9bcc0f562df401aac051912a88658feb6f1d43344}
\newcommand{\HashLeanTrialDivisionLimit}{48212906b2df9199981ff04a4a298ceaba35bf1c9205be786f2290b5a9f422ad}
\newcommand{\HashLemmaCompositeHasPrimeFactorLeSqrt}{a99f4d1a37cf56b8abf2cd6ea57e9a2e3ff094d41c514059d3a216d3c024a87d}
\newcommand{\HashLemmaCoprimeImpliesMultiplicative}{43dc1ca8a2341b479a646e0f777e803822f2f83d74f1047e362b1b7d9bd73a54}
\newcommand{\HashLemmaCoprimeImpliesMultiplicativeNonlinear}{3f83199ee284d4372f8742bee07b3240d22b0a0d28299f473a3220576b7645a8}
\newcommand{\HashLemmaDisjointByConstruction}{7c9915242c5363d0d746cf738fcc8d5a645e83a0a97b4c81cb5803a5606c497e}
\newcommand{\HashLemmaDivisibilityBounds}{21db37ec4b3a1fc81447e1a6d7662b5491f77c0007e2572f7aab0893f3094fad}
\newcommand{\HashLemmaDivisibilityTransitive}{2ec5734be5cfa0240e2da57f8902667eb1297af8da8ff31c4fc0067c36ace4e8}
\newcommand{\HashLemmaFSquaredGtNMinusOne}{0bb50154bb815615185f7e10607fd00939b27471e43f6e0b2fad3fc358e5755b}
\newcommand{\HashLemmaFermatLittleTheorem}{a3b43a62b0532015c72eb3b25babbcbfc3b87f8eed96be75527c8e2c07bd1dc1}
\newcommand{\HashLemmaModpowAddMul}{5c061a58da2b5632f44f2ce4760521ec9766913e7230e8bbda3f28c04d4d49f4}
\newcommand{\HashLemmaModpowModDivisibility}{7d8df134ae16613cfc7b10a8cedb6fe351d204f6f8fabf235d4ef651cf6a6404}
\newcommand{\HashLemmaOrderExists}{83d837d2873f69f81afc7c8e5ba2accfa6d1b22d221889de7ad50de1bc89fa8a}
\newcommand{\HashLemmaOrderLePMinusOne}{afa07507ab9c53b6308ef849104ebc11135654d35f4415f472ded94328eee5ac}
\newcommand{\HashLemmaOrderPrimeFactor}{dbc5c22b9ad8fa4b277f45559460f48728f30b27918519d39e37aa1603107b9e}
\newcommand{\HashLemmaPocklingtonCertificate}{4c18755e2796218d2dc167b161ee2be96223bcdbf2b75c3e1c7ab7274105cb21}
\newcommand{\HashLemmaSigmaMultiplicative}{e5a0eb59ed740c0608f7464e8076cf9a6938711b4532d7b5db353c65c86e7bf8}
\newcommand{\HashLemmaSmallestFactorIsPrime}{eef56c28181f254140a52a6d139507ab9712c59c5330df577a98f871243af286}
\newcommand{\HashLemmaSquareComparisonContradiction}{001fe983f651d95c3d53b524e3fbb19e1d697fc8fe0ec4db056a86dbc3a80f2b}
\newcommand{\HashModpowSpec}{57bd9cccef197081a718cbe176d531117f970030cf1f0b63b315e7ee458ca6a9}
\newcommand{\HashNatToUFiveOneTwo}{ab0f53d68438140908c8ad64a6963494608dd4eb6a7c6736ac21b37f80d1d003}
\newcommand{\HashPassesRaycastSieveSpec}{3383c2d72949bc8cc1a32bd399ddaf5e51dff95c539e7a29da11e308e1617182}
\newcommand{\HashPocklingtonSpec}{72a6261d17dae37711c1877870920465513addc9e75c8c7fb14937805330f3ef}
\newcommand{\HashPowTwoSixFour}{f5af310449d93b35707c205bbb1f8c9e0ef77725cdb4dae9fe1d67d6fcea4adb}
\newcommand{\HashPowSpecVal}{4083ae88a2a2d218fa5119510b45ac300b7c4a0d1a5e934305546877c9eab2fd}
\newcommand{\HashPrasadSunithaBoundSatisfied}{8e55cd49edcfc1f202b00f41b27d896612482209565483df7f3c239072849bbe}
\newcommand{\HashPrimeFactors}{7aaec2c92cf85360eef84f62efea4c5853ead314001cb136f3f3434976666e91}
\newcommand{\HashProofVerifyZsigmondyPreconditions}{176c508e942613bfb5b232045b786a9c729703342f7e7cd1ba2d445a20a4ffa2}
\newcommand{\HashProveBaselineMinPrimeFactorsEquivalence}{533d6e464fa3d8148b3bf3beeaac24b6401cba8df3470206b25236506258e0da}
\newcommand{\HashProveCombinedBounds}{c620b68e324805d0dcde60cba5e2d25904cb41f2595beaaeb6386e4d6a40137e}
\newcommand{\HashProveConjecturalActiveEquivalence}{c9c71c6f464007e269f2deed65ed1e98d278f81d68a9430f500f3cebdc346dfd}
\newcommand{\HashProveConjecturalMaxLogOneZeroCeilingEquivalence}{f2146249b5573c6438e3018e67479febfaf79e3a27478209ca7f0b03f85ae138}
\newcommand{\HashProveCrtModulusProductEquivalence}{a45f4d6de241034e881a75e1139798998ea261a1e9cc115c9e1bca04a3757933}
\newcommand{\HashProveDivFiveCoprimeThreeBoundEquivalence}{4ad066861c3a28a49d8f22701e03e426d6333285c3a6231ed181b44e223f435a}
\newcommand{\HashProveDivFiveCoprimeThreeCombinedEquivalence}{499016bc6febaec0360930d79bbc1988b55b4a830b2b124b4421f6ca3e766529}
\newcommand{\HashProveEulerCeilingDenEquivalence}{c750bc0fa7107b5de2a908f112da90261c50570a0e1098ce91afa38345dff0ec}
\newcommand{\HashProveEulerCeilingNumEquivalence}{420ed686357773e11c1ebd88e3581fdde67eb2fd45ee4287aaca5b5a0de9b76c}
\newcommand{\HashProveMaxExponentEquivalence}{19e5f503000b8c640e67a7d3b6b208c6af0eb0a46e201312242162d8944793e7}
\newcommand{\HashProveOverflowThresholdDenEquivalence}{0c9a5b13eb2a719d7adac8adedbe4818c50c47ece7d2e695798dc523d9665e1e}
\newcommand{\HashProveOverflowThresholdNumEquivalence}{629295535605e1d444a12dd4816f2ab33d882f087fb6563789af647f94efc9e8}
\newcommand{\HashProvePollardRhoBatchSizeEquivalence}{cc06e97aeffdd6b05a0ca5da2a1d9b4c17523acb48ccdf6daf3e05e42c4ff064}
\newcommand{\HashProvePollardRhoIterationLimitEquivalence}{7b2346c36d4b490efc5354e58521a50ea658b4390073d35a2f459914e4544f5c}
\newcommand{\HashProvePrasadSunithaBoundEquivalence}{7ad9a4710c848b4444823d8322306e00ce525249116927f82862db6742af2296}
\newcommand{\HashProvePrasadSunithaCombinedEquivalence}{8c5b9a71d31ad70843b4616caa438c671c8c9628ddfbad9398a71918e16912a6}
\newcommand{\HashProvePrefixStopThresholdEquivalence}{956fb011b8e9133bbdbfaae54821f00ee6541a7b9005ff9e87ceaae87cdfbd93}
\newcommand{\HashProvePrimeSplitThresholdEquivalence}{3d35da3451eb6a39a71cfc307d38c12d76bea9429cc4a43a80976cd90cc7f3ae}
\newcommand{\HashProveRaycastChunkSizeEquivalence}{94bdcfac7dd611886967ee4ce1f5b57b4be705505a561d3b82618a4e453e6d8e}
\newcommand{\HashProveRaycastGpuThresholdEquivalence}{f8a5a1468d505de60215a289ed7e8d78480f70ebffb4f4bf09e55121fdb991a5}
\newcommand{\HashProveSieveLimitEquivalence}{1cc13a1fa752a539fc523e7526d42c2afebd6c21bb86a727cd7f651777bcea6c}
\newcommand{\HashProveTargetMaxLogOneZeroEquivalence}{b4932d34fa506bc6cc57750c9b392edb7717d1e74df198e8035bbf0befaf9d32}
\newcommand{\HashProveTargetMinLogOneZeroEquivalence}{3e6375b921b467b0f1ea20cec033c0e55934f10a8fa5dab4d39b135f9488449f}
\newcommand{\HashProveTrialDivisionLimitEquivalence}{11125a55ede2e273ea2cfa64f2362c2d2118edd4980835117ac475fc052b0b48}
\newcommand{\HashRnsFiveOneTwoGcdSpec}{326bf8dad421fb30d03432967c2fa2e50077f926b6b04d96fed84e8ac39fede6}
\newcommand{\HashRnsFiveOneTwoMontMulSpec}{8bf52d4e2e0e21a29f600c907229fd2c9ec76c010f5524ea15ce97c3c5ed48c7}
\newcommand{\HashRnsFiveOneTwoValidBounds}{bfe913f82b1159919d12287b03b1ebd1c7816b606e1e3c4b892fbb217d90f5e7}
\newcommand{\HashScaleBoundCeil}{209ae501b1cf4f71b3b81e6a00cabcdc8e5fc1472d8c16d80f401d3e49c340bf}
\newcommand{\HashScaleBoundSpec}{def035629b1ef362535899872bea3405160c6b7f3836dd9f4b267f7634c12d03}
\newcommand{\HashScreenModEight}{99e0f1eedd8c0a5ed8de911f482f0bf02427ecfe29ad84391aaf2284b2abede1}
\newcommand{\HashSigmaSpec}{dce126e7cfef78866bffcb84dfeb0f82330850c0ef177a2f3166540aa6405f9b}
\newcommand{\HashSigmaSpecVal}{d9a7fb2dbe4e4d37f1c9752a6aac71abab205b37829db533cfa0ce8355ac8ae8}
\newcommand{\HashTryScaleBoundCeil}{b557cc466b504830ce067946b3541e2741a904d4ae14659ec620f39ae02e5020}
\newcommand{\HashUFiveOneTwoToNat}{2515c2db85b32ffba43960532624bc0c6a371d52d528926e06de6e8da935a77d}
\newcommand{\HashVerifiedAllocUFiveOneTwo}{725a44210587084ad50473e06994a60da9cb63463c4fd92a9e0a5fd93180789a}
\newcommand{\HashVerifiedFreeUFiveOneTwo}{d53f5bca387920db013e709a9cdb5f267318857ceb57505ca1755e87b0927ab3}
\newcommand{\HashVerifiedGetUFiveOneTwo}{3d0b2b77ca741a04cfbf52fffe7492f1cdd4cf18afbdb8d12e38c8f672b5b6d1}
\newcommand{\HashVerifiedPassesRaycastSieve}{98d0f8bdeafdb39be4a396c79598aad9adc84bebf35f50b70916cf1768740764}
\newcommand{\HashVerifiedUalbfComputeSigma}{ebe1a16d20de6e44aa75290dc6cc1033f57342d5adc05cec6f28a59a11ebe58b}
\newcommand{\HashVerifiedUalbfCyclotomicEval}{9aa2ae305bcc73065cb6dfd9b5270a8452cb14e7d69d0a91d1ada36322ef9bc1}
\newcommand{\HashVerifyPrasadSunitha}{8a144688b3e77868b23afa741f728d0f1646ab2734eb993fd0a8ecbea9394251}
\newcommand{\HashVerifyStarvationPruning}{3fe33985a5f6827b419338c23cc10bb0e50ea2d8bb4dc4ce39c1ccc6a8e95e9c}
\newcommand{\HashZsigmondyPreconditionsSatisfied}{f105d37de697516ec3aa368672cb12427ad608ff9da21baa85cf61f17c35709f}


\maketitle

%% =================================================================
%%  Abstract
%% =================================================================
\begin{abstract}
This paper presents the Unified Algebraic-Lattice Bipartition Framework
(UALBF), a formal-verification and computational
approach to the quasiperfect number (QPN) problem. Using Lean~4,
we mechanically verify that every QPN is an odd perfect square
(a result of Cattaneo) and
establish a modulo-8 obstruction on prime factors of $\sigma(N)$.
We prove the abundancy index $H(N)$ is bounded by the Euler
totient ratio $N/\varphi(N)$, decomposed into $H(N)$ times a correction
factor $C < 1.022$. This leverages pure-$\mathbb{Q}$ telescoping-sum identities
(the 61/60 tail bound) without real-analysis imports. The chain
bounds $N/\varphi(N) < 2.0442$ for QPNs coprime to~15.

We further formalise the Prasad--Sunitha bound $\omega(N)\ge 15$ for
$\gcd(N,15)=1$, and a fully decomposed structural proof of Zsigmondy's
theorem via seven modular sub-lemmas, isolating analytic ingredients
into independently verified theorems.

These lemmas are used by a parallel Rust search engine, which we
describe as an engineering prototype. Its coverage is incomplete:
it only considers primes below 250{,}000 with exponents $2e \le 8$,
one pruning rule treats primes that a Lean theorem shows divide
$\sigma(N)$ as if they had to divide~$N$, several pruning bounds only account for
components in its finite list, and its exact ray-casting stage did
not run. The Lean library proves individual pruning lemmas; nothing
proves that the search as a whole is complete, and this paper does
not establish any new lower bound on quasiperfect numbers. The
strongest bound we are aware of is $N > 10^{45}$ (Alekseyev, 2026).
\end{abstract}


%% =================================================================
%%  Section files
%% =================================================================
%% =================================================================
%%  1.  Introduction
%% =================================================================
\section{Introduction}
\label{sec:intro}

A quasiperfect number (QPN) is a positive integer~$N$ whose sum of
divisors satisfies $\sigma(N) = 2N + 1$.  No such number is known.
Cattaneo (1951) showed that any QPN must be an odd perfect
square~\cite{cattaneo1951}.  Hagis and Cohen (1982) showed that a QPN
satisfies $N > 10^{35}$ and has at least seven distinct prime
factors~\cite{hagis1982}.  Prasad and Sunitha (2017) showed that a QPN
with $\gcd(N,15)=1$ has at least fifteen distinct prime
factors~\cite{prasad_sunitha}.  More recently, Alekseyev (2026) showed
that there is no QPN below $10^{45}$~\cite{alekseyev2026}, and
Toyohara, Tao and Yao (2026) claim that every QPN has at least eight
distinct prime factors~\cite{toyohara_tao_yao2026}; the latter result
should be treated as conditional until their enumeration code is made
public.  Pushing these bounds further by computation is difficult
because the number of possible prime factorisations grows
combinatorially.

\paragraph{The verification gap.}
The central obstacle confronting modern computational attacks on the QPN
problem is not raw arithmetic speed: it is the \emph{verification gap}.
Purely mathematical approaches struggle to eliminate vast swaths of
candidate numbers without concrete enumeration.  Purely computational
searches, meanwhile, are prone to unverified optimisations and integer
overflow bugs that might silently skip a valid QPN\@.  Translating
formally proved bounds into executable code introduces a fresh class of
risks---floating-point inaccuracies, off-by-one errors in sieve
implementations, and race conditions in parallel search---that can
compromise mathematical soundness without any visible failure.

\paragraph{Our contribution.}
The Unified Algebraic-Lattice Bipartition Framework (UALBF) attacks the
verification gap head-on.  Instead of presenting mathematics and
computation as separate concerns, the UALBF \emph{fuses} them: the Rust
search engine calls compiled Lean~4 binaries through a C-level Foreign
Function Interface (FFI)\@.  Critical hot-path operations---128-bit
modular inverses for ray-cast targeting and exact $\sigma(p^{2e})$
evaluations---are dispatched into Lean code whose correctness is
established by formal bridge theorems: \texttt{computeSigmaNat\_eq\_sigma}
proves that the FFI's $\sigma$ computation equals Mathlib's
sum-of-divisors function, and \texttt{extGcdAux\_bezout} establishes
B\'ezout's identity for the extended GCD underlying every modular inverse.
Runtime overflow is prevented by \texttt{\_ok} sentinel exports
(\texttt{ualbf\_compute\_sigma\_ok}, \texttt{ualbf\_mod\_inverse\_ok})
that verify results fit within 128 bits before the Rust engine reads the
truncated word halves.  The result is a system that executes
\emph{mathematically certified, C-linked binaries} inside a high-speed,
lock-free search engine.

\smallskip
In this paper, we present the following core contributions:

\begin{itemize}
  \item \textbf{Formalized Parity and Divisibility Lemmas:}
    We provide mechanically verified proofs in Lean~4 that a QPN cannot
    be a double square, and establish the foundational $\sigma$ parity
    and exact valuation lemmas (\cref{sec:math_foundations}).

  \item \textbf{Abundancy Index and Totient Ratio Analysis:}
    We formally verify that the abundancy index
    $H(N) = \sigma(N)/N$ is strictly bounded by the Euler totient ratio
    $N/\varphi(N)$, decompose this ratio into $H(N) \times C$ where $C$
    is a correction factor product, and prove $C < 1.022$ for QPNs
    coprime to~15, yielding $N/\varphi(N) < 2.0442$
    (\cref{subsec:abundancy_euler_ceiling,subsec:totient_decomposition}).

  \item \textbf{Cyclotomic Factorisation and Zsigmondy Obstructions:}
    We formalise the cyclotomic polynomial decomposition of
    $\sigma(p^{2e})$ and decompose the proof of Zsigmondy's theorem
    into seven modular sub-lemmas, establishing
    that high exponents spawn primitive prime divisors triggering the
    modulo-8 obstruction (\cref{subsec:cyclotomic_zsigmondy}).
    All sub-lemmas are now formally verified in Lean~4/Mathlib.
    The decomposition exposes every analytic ingredient
    as a separate lemma, creating a 100\% mechanically verified
    cyclotomic decomposition with zero gaps.

  \item \textbf{Standalone Correction Factor Module:}
    We provide a pure-$\mathbb{Q}$ telescoping-sum framework in
    \texttt{Pure/RationalBounds.lean} (the 36/35 pure identity machinery), 
    which serves as a direct mathematical dependency to bound the product 
    tail for the unified 2.0442 Euler ceiling theorem, requiring no real-analysis imports
    (\cref{subsec:correction_factor_standalone}).

  \item \textbf{Bipartition Congruence Lemma:}
    We prove the contrapositive
    \texttt{no\_\allowbreak{}solution\_\allowbreak{}no\_\allowbreak{}qpn}:
    if a suffix $N_R$ fails the congruence
    $N_R \cdot 2N_L \equiv -1 \pmod{\sigma(N_L)}$, then
    $N_L N_R$ is not quasiperfect (\cref{subsec:bipartition}).  This is
    a lemma about a single prefix--suffix pair; it does not show that
    the Rust engine enumerates every admissible suffix.

  \item \textbf{Theoretical Framework for Deeper Horizons: Exact Valuation Ray-Cast Strategy:}
    We introduce a ray-cast algorithm in Rust whose \emph{targeting
    step}---computing the unique residue class of $N_R$ modulo
    $\sigma(N_L)$ via a single modular inverse---runs in
    $\mathcal{O}(1)$ time. Candidates along the resulting arithmetic
    progression are then enumerated.  This stage did not run in the
    search reported here (\cref{sec:computational_methodology}).

  \item \textbf{Prototype Search Engine:}
    We describe a prototype Lean/Rust search engine and list the
    reasons its coverage is incomplete (\cref{sec:results}).  The search
    does not establish a lower bound on quasiperfect numbers.
\end{itemize}


%% -------------------------------------------------------------------
\subsection{The Quasiperfect Number Problem}
\label{sec:qpn_problem}

For a positive integer~$N$, the \emph{sum-of-divisors function} is
$\sigma(N) = \sum_{d \mid N} d$.  The \emph{abundancy index} of~$N$ is
the ratio $\sigma(N)/N$.  A \emph{perfect number} satisfies
$\sigma(N)=2N$; analogously, a \emph{quasiperfect number} (QPN) is a
positive integer satisfying
%
\begin{equation}\label{eq:qpn-def}
  \sigma(N) \;=\; 2N + 1.
\end{equation}

\begin{definition}[Quasiperfect Number]\label{def:qpn}
  A positive integer~$N$ is \textbf{quasiperfect} if $N > 0$ and
  $\sigma(N) = 2N + 1$.  In our Lean~4 formalisation this is encoded
  directly:
  \begin{center}
    \leaninline{def IsQuasiperfect (n : ℕ) : Prop := n > 0 ∧ sigma n = 2 * n + 1}.
  \end{center}
\end{definition}

No quasiperfect number has ever been found.  Cattaneo~\cite{cattaneo1951}
showed that any QPN must be an odd perfect square, and Hagis and
Cohen~\cite{hagis1982} established that $N > 10^{35}$ with at least
seven distinct prime factors.  Because $2N+1$ is manifestly odd, one
obtains immediately that $\sigma(N)$ is odd for any QPN; the
consequences of this parity constraint are the starting point for the
analysis below.

\paragraph{Setup.}
The project setup explicitly relies on our foundational dependencies: Lean 4~\cite{moura2021}, Mathlib~\cite{mathlib2020}, Rayon~\cite{rayon2024}, bnum~\cite{bnum2024}, and ed25519-dalek~\cite{ed25519dalek2024}.

\input{sections/02_math_and_formalization}
%% =================================================================
%%  3.  The Unified Algebraic-Lattice Bipartition Framework
%% =================================================================
\section{The Unified Algebraic-Lattice Bipartition Framework}
\label{sec:ualbf}

By \cref{thm:qpn-odd-sq}, every QPN is an odd perfect square
$N = \prod_{i} p_i^{2e_i}$
with each $p_i$ an odd prime.  The key idea of the UALBF is to
\emph{bipartition} this factorisation into a \emph{prefix} $N_L$ and a
\emph{suffix} $N_R$, chosen so that all arithmetic constraints induced by the
QPN equation decompose cleanly across the split.  The prefix is constructed
incrementally by the Rust engine's depth-first search, while the suffix's
residue class modulo $\sigma(N_L)$ is determined by a constant-time
modular inverse; candidates in that class are then enumerated and sieved.

%% -------------------------------------------------------------------
\subsection{Theoretical Frameworks for Deeper Horizons: Bipartitioning the Search Space}
\label{subsec:bipartition}

\emph{(Note: the ray-casting stage described in this section and in
\cref{subsec:ray_casting} did not run in the search reported in
\cref{sec:results}.)}

\begin{definition}[QPN Bipartition]\label{def:bipartition}
A \textbf{QPN bipartition} of a quasiperfect number $N$ is a triple
$(N, N_L, N_R)$ satisfying:
\begin{enumerate}
  \item $N_L > 0$, $N_R > 0$;
  \item $N = N_L \cdot N_R$;
  \item $\gcd(N_L, N_R) = 1$.
\end{enumerate}
Any partition of the prime set of $N$ into two disjoint subsets induces
a valid bipartition; the Rust engine's DFS assigns the first $k$ primes
to $N_L$ and all remaining primes to $N_R$.
\end{definition}

\paragraph{Multiplicativity over the bipartition.}
Because $\gcd(N_L, N_R) = 1$, the sum-of-divisors function factors:
\begin{equation}\label{eq:sigma-bipartition}
  \sigma(N) \;=\; \sigma(N_L)\,\sigma(N_R).
\end{equation}
Substituting the QPN equation $\sigma(N) = 2N + 1 = 2 N_L N_R + 1$ yields
the \emph{fundamental bipartition identity}:
\begin{equation}\label{eq:bipartition-identity}
  \sigma(N_L)\,\sigma(N_R) \;=\; 2\,N_L\,N_R + 1.
\end{equation}

\paragraph{Coprimality of prefix and its divisor sum.}
The identity~\eqref{eq:bipartition-identity} has an important structural
consequence:

\begin{theorem}[Prefix--Sigma Coprimality]\label{thm:prefix-coprime}
For any QPN bipartition $(N, N_L, N_R)$,
\[
  \gcd\bigl(N_L,\;\sigma(N_L)\bigr) \;=\; 1.
\]
\end{theorem}

\begin{proof}
Let $d = \gcd(N_L, \sigma(N_L))$.  We show $d \mid 1$.

Since $d \mid N_L$, we have $d \mid N_L \cdot N_R$, hence
$d \mid 2\,N_L\,N_R$.
Since $d \mid \sigma(N_L)$ and $\sigma(N_L)$ is a factor of
$\sigma(N_L)\cdot\sigma(N_R)$, we have $d \mid \sigma(N_L)\cdot\sigma(N_R)$.
By the bipartition identity~\eqref{eq:bipartition-identity},
$\sigma(N_L)\cdot\sigma(N_R) = 2N_LN_R + 1$.
So $d$ divides both $\sigma(N_L)\cdot\sigma(N_R)$ and
$2N_LN_R$, hence $d$ divides their difference:
$d \mid (2N_LN_R+1) - 2N_LN_R = 1$.
Therefore $d = 1$.
\end{proof}

\cref{thm:prefix-coprime} is operationally critical: it guarantees that the
modular inverse of $2 N_L$ modulo $\sigma(N_L)$ always exists, so the Rust
engine's ray-cast subroutine can never encounter a division-by-zero panic.

\paragraph{The AMBS suffix target.}
Casting~\eqref{eq:bipartition-identity} into the ring
$\mathbb{Z}/\sigma(N_L)\mathbb{Z}$, the left-hand side vanishes (since
$\sigma(N_L) \equiv 0$), giving
\[
  0 \;=\; 2\,N_L\,N_R + 1
  \qquad\text{in } \mathbb{Z}/\sigma(N_L)\mathbb{Z}.
\]
Rearranging:
\begin{equation}\label{eq:ambs-target}
  N_R \cdot (2\,N_L) \;\equiv\; -1 \pmod{\sigma(N_L)}.
\end{equation}
Because $\gcd(2 N_L, \sigma(N_L)) = 1$ (by \cref{thm:prefix-coprime} and
the fact that $\sigma(N_L)$ is odd for a QPN), the value of $N_R$ modulo
$\sigma(N_L)$ is uniquely determined.  This is the \emph{Algebraic-Modular
Bipartition Sieve} (AMBS) target: given a prefix $N_L$, the engine computes
the unique residue class of $N_R$ in $\mathcal{O}(1)$ time via a single
modular inverse, then enumerates all representatives
$z \equiv r \pmod{\sigma(N_L)}$ in the interval $[z_{\min},\, z_{\max}]$
in $\mathcal{O}(z_{\max}/\sigma(N_L))$ time, applying the mod-8 sieve
and coprimality checks to each candidate.

The AMBS target equation~\eqref{eq:ambs-target} also admits a clean
contrapositive:

\begin{theorem}[Bipartition Congruence Lemma]\label{thm:no-sol-no-qpn}
If $N_R \cdot (2\,N_L) \not\equiv -1 \pmod{\sigma(N_L)}$ for all
valid suffix candidates, then $N = N_L \cdot N_R$ is not quasiperfect.
\end{theorem}

\begin{proof}
Suppose for contradiction $N$ is quasiperfect.
By~\eqref{eq:ambs-target}, any such $N$ must satisfy
$N_R \cdot (2N_L) \equiv -1 \pmod{\sigma(N_L)}$.
But every valid suffix candidate fails this congruence by hypothesis,
so the QPN hypothesis implies a contradiction.
\end{proof}

In the Lean formalisation, \cref{thm:no-sol-no-qpn} is
\texttt{no\_\allowbreak{}solution\_\allowbreak{}no\_\allowbreak{}qpn}
in \texttt{Bipartition.lean}.  It justifies discarding any suffix that
fails the congruence~\eqref{eq:ambs-target}.  It says nothing about
whether the Rust engine actually enumerates every admissible suffix:
that depends on the search bounds, on the finite list of components
the engine considers, and on the correctness of the enumeration code,
none of which is formally verified.  The lemma therefore does not turn
a ray-cast run into a non-existence proof for a prefix.



This component heavily relies on our core software dependencies, specifically Lean 4~\cite{moura2021}, Mathlib~\cite{mathlib2020}, Rayon~\cite{rayon2024}, bnum~\cite{bnum2024}, and ed25519-dalek~\cite{ed25519dalek2024}.

%% =================================================================
%%  4.  The Verified Computational Engine
%% =================================================================
\section{The Verified Computational Engine}
\label{sec:computational_methodology}

%% -------------------------------------------------------------------
\subsection{Closing the Verification Gap via Lean FFI}
\label{subsec:lean_ffi}

A fundamental challenge in computational number theory is the
``verification gap'': translating abstract formal proofs into
executable, performant code without introducing unverified
implementation errors (e.g., floating-point inaccuracies or integer
overflows).  To close this gap, the UALBF architecture relies on a
compiled Lean Foreign Function Interface (FFI) bridge
(\texttt{FFI.lean} and \texttt{lean\_\allowbreak{}ffi.rs}).

Instead of re-implementing validated mathematical functions in Rust, we
author executable definitions directly in Lean~4 (e.g.,
\texttt{computeSigmaNat} and \texttt{modInverse}).
These functions are tagged with \texttt{@[export]} and compiled into a
static C-core library that the Rust orchestrator binds to via FFI\@.

\paragraph{Formal bridge theorems.}
Because Lean erases \texttt{Prop} proofs at runtime, the executable
definitions in \texttt{FFI.lean} are \emph{not} verified by
construction alone---they are standalone computational algorithms.
To close the epistemological gap, we provide formal bridge theorems
that prove equivalence between the FFI executables and the mathematical
definitions used in the proof library:

\begin{itemize}
  \item \textbf{\texttt{computeSigmaNat\_eq\_sigma}}: Proves that
    \texttt{computeSigmaNat~p~e} equals
    $\sigma(p^e) = \sum_{d \mid p^e} d$ for prime~$p$, composing
    Mathlib's \texttt{sum\_divisors\_prime\_pow} with the geometric sum
    identity \texttt{nat\_geom\_sum}.
  \item \textbf{\texttt{extGcdAux\_bezout}}: Proves B\'ezout's identity
    $a x + b y = g$ for the extended GCD, establishing the core invariant
    underlying the modular inverse computation.
  \item \textbf{\texttt{modInverse\_spec}}: Proves that when
    \texttt{modInverse a m} returns \texttt{some~v}, we have
    $(a \cdot v) \equiv 1 \pmod{m}$, closing the correctness gap between
    the FFI's extended-GCD computation and genuine modular invertibility.
    This theorem is fully proven---no \texttt{sorry}s remain.
\end{itemize}

\paragraph{Overflow guards.}
To prevent silent truncation when splitting arbitrary-precision Lean
\texttt{Nat} values into 64-bit word pairs, every exported function
pairs its \texttt{\_lo}/\texttt{\_hi} word outputs with a dedicated
\texttt{\_ok} sentinel:
\begin{itemize}
  \item \textbf{\texttt{ualbf\_compute\_sigma\_ok}}: Returns~1 iff
    $\sigma(p^e) < 2^{128}$; the Rust engine treats a~0 return as a
    \texttt{None} result and skips the component.
  \item \textbf{\texttt{ualbf\_mod\_inverse\_ok}}: Returns~1 iff
    $\gcd(a, m) = 1$, i.e.\ the inverse exists and fits in the
    128-bit word pair; a~0 return signals that the ray-cast targeting
    step is inapplicable.
\end{itemize}
The Rust engine reads \texttt{\_lo}/\texttt{\_hi} only when the
corresponding \texttt{\_ok} flag is non-zero, guaranteeing that
no modulo-truncated garbage enters the search arithmetic, including
if the search bounds are raised. (See \cref{app:verification} for an exhaustive technical appendix detailing the low-level execution invariants, specifically the Pocklington primality testing and RNS512 arithmetic models used within the computational engine.)

This architecture integrates the FFI bridge as an executable boundary atop our
stratified four-layer Lean formalisation (\cref{sec:math_foundations}): the executable
definitions in \texttt{FFI.lean} depend on the Engine and QPN tiers for their
correctness guarantees, but are compiled into standalone C-callable
binaries that the Rust engine invokes without any Lean runtime overhead.


%% -------------------------------------------------------------------
\subsection{Phase~1: The Global Annihilation Sieve}
\label{subsec:global_sieve}

The UALBF computational engine is driven by a Lean-managed loop
(\texttt{ualbf\_\allowbreak{}dfs\_\allowbreak{}loop\_\allowbreak{}impl}) that fuses the generation of candidate prefixes with
immediate algebraic validation.  The engine traverses an array of valid
prime power components $p^{2e}$ using a Lean-orchestrated Depth-First Search (DFS)
architecture. Lean retains absolute control over the search space flow, making deterministic
state transitions via FFI calls (push, backtrack, evaluate). To eliminate memory allocation bottlenecks,
the search operates on a sequential, zero-allocation push/pop stack model.
Throughout the DFS, active branches are subjected to
aggressive mathematical sieves, dynamically tracking accumulated metrics
such as the running abundancy ratio
$\prod \sigma(p_i^{2e_i})/p_i^{2e_i}$ and enforcing strict divisibility
constraints to prune unviable subsets early.

Before any DFS traversal begins, the engine precomputes the universe of
valid prime power components via the \emph{Global Annihilation Sieve}
(\texttt{phase1\_\allowbreak{}global\_\allowbreak{}annihilation\_\allowbreak{}sieve}).
For each odd prime~$p$ up to 250{,}000 and exponent
$2e \in \{2, 4, 6, 8\}$, the engine computes $\sigma(p^{2e})$ using the
formally verified Lean FFI (\texttt{compute\_\allowbreak{}sigma}) and
factors it using cyclotomic polynomial decomposition.
Prime powers outside this range (primes $p \ge 250{,}000$, or
exponents $2e > 8$) are never generated, so the search does not
consider any $N$ with such a component; this is a limitation of
coverage, not something the Lean library justifies.

\paragraph{Cyclotomic decomposition for efficient factorisation.}
Rather than factoring the full value $\sigma(p^{2e})$ directly---which
can exceed $2^{128}$---the engine exploits the cyclotomic expansion of
\cref{lem:cyclotomic-expansion}.  The routine
\texttt{factor\_\allowbreak{}sigma\_\allowbreak{}cyclotomic} computes
each cyclotomic factor $\Phi_d(p)$ for $d \mid (2e+1)$, $d > 1$, via
the M\"obius inversion identity
$\Phi_d(x) = \prod_{k \mid d} (x^k - 1)^{\mu(d/k)}$.  Each
$\Phi_d(p)$ is typically much smaller than the full $\sigma$ value,
enabling the Elliptic Curve Method (ECM) factoriser (via the
\texttt{prime\_\allowbreak{}factorization} crate) to complete
efficiently.  When a cyclotomic evaluation overflows $\mathtt{u128}$,
the engine falls back to factoring the full $\sigma$ value computed via
the Lean backend.

\paragraph{Mod-8 annihilation.}
Every prime factor~$q$ of $\sigma(p^{2e})$ is checked against the
modulo-8 obstruction (\cref{thm:mod8}): if $q \equiv 5$ or
$7 \pmod 8$, the component $p^{2e}$ is discarded.  By
\cref{thm:sieve-sound}, this pruning is sound---no legitimate QPN
factor is lost.

The surviving components are sorted by abundance ratio
$\sigma(p^{2e})/p^{2e}$ in descending order, placing the most
productive factors first for the DFS.

%% -------------------------------------------------------------------
\subsection{Depth-First Search for Prefix Construction}
\label{subsec:dfs_prefix}

The engine constructs candidate prefixes via a Lean-managed DFS over the
sieved components, mapping the formalized exact valuation
constraints directly onto the computational state.

\paragraph{Lean orchestration and FFI state management.}
Rather than delegating control flow to Rust, Lean orchestrates the primary
search loop (\texttt{ualbf\_\allowbreak{}dfs\_\allowbreak{}loop\_\allowbreak{}impl}). Lean retains absolute control over the search space
flow, executing deterministic state transitions via explicit FFI bridge protocols:
\begin{itemize}
  \item \textbf{Push:} Lean dictates when a valid component extends the current prefix, invoking FFI calls to mutate the underlying Rust state stack.
  \item \textbf{Evaluate:} At each node, Lean requests pruning evaluations from the Rust engine to dynamically calculate metrics such as abundancy limits.
  \item \textbf{Backtrack (Pop):} When a branch is exhausted or pruned, Lean explicitly signals the FFI to pop the state, ensuring the formal proof checker maintains a continuous, verified chain of the traversal history.
\end{itemize}

\paragraph{Breadth-level parallelism.}
The combinatorial explosion of possible QPN prime factors necessitates parallel evaluation without surrendering flow control. Rust's parallelism, via the Rayon framework, is strictly limited to breadth-level component selection. At any given depth layer dictated by Lean, Rayon utilizes work-stealing to concurrently evaluate and filter the valid siblings (components) for the next potential push. Rust does not control parallel depth search paths, ensuring that Lean's sequential verification model remains the undisputed source of truth.

\paragraph{Running abundancy tracking.}
Each \texttt{Prefix} struct carries a
\texttt{current\_\allowbreak{}abundancy} field---the running product
$\prod_{i} \sigma(p_i^{2e_i})/p_i^{2e_i}$ over all components assigned
to the prefix.  This field is updated incrementally (multiply on push,
restore on pop) and drives several pruning mechanisms:

\begin{itemize}
  \item \textbf{Overflow Kill:}
    If $\texttt{current\_\allowbreak{}abundancy} > 2.000001$, the
    prefix is immediately discarded.  Since
    $H(N) = 2 + 1/N \approx 2$ for large~$N$, exceeding this
    threshold indicates an impossible abundancy overrun.

  \item \textbf{Suffix Starvation Kill:}
    Before each DFS extension, the engine checks whether
    $\texttt{current\_\allowbreak{}abundancy}
     \times \texttt{suffix\_\allowbreak{}abundance}[i]$
    can reach~2, where $\texttt{suffix\_\allowbreak{}abundance}[i]$ is
    a precomputed array storing the maximum achievable abundancy product
    from the remaining components.  If not, the branch is pruned
    (\cref{thm:starvation}).  The precomputed array only accounts for
    components in the engine's finite list, so this pruning is sound
    only relative to that list, not for arbitrary~$N$.

  \item \textbf{Ray-Cast Space Exhaustion:}
    When a prefix magnitude reaches the stopping threshold, the engine is designed to hand the remaining subspace to the exact ray-casting subroutine and stop extending that prefix.  The branch is then only as complete as the ray-cast enumeration, which is not formally verified and did not run in the search reported in \cref{sec:results}.
\end{itemize}

%% -------------------------------------------------------------------
\subsection{Theoretical Frameworks for Deeper Horizons: Orchestration and Ray-Casting}
\label{subsec:ray_casting}

\emph{(Note: the ray-casting stage described in this section did not
run in the search reported in \cref{sec:results}.)}

The ray-cast strategy serves as the core enumeration engine of the
UALBF\@.  Given a prefix~$N_L$, the algorithm must find valid suffix
candidates $N_R = z^2$.  From the AMBS target
(\cref{eq:ambs-target}), we have
$N_R \equiv -(2N_L)^{-1} \pmod{\sigma(N_L)}$.  Let $X_L$ be this
target residue.  The targeting step, executed via a fast 128-bit
modular inverse function
(\texttt{mod\_\allowbreak{}inverse\_\allowbreak{}128}) validated through
the Lean FFI, establishes that
$z^2 \equiv X_L \pmod{\sigma(N_L)}$.

To enumerate candidate roots~$z$, the engine employs the
Tonelli--Shanks quadratic root extraction algorithm, generalised for
composite moduli via the Chinese Remainder Theorem in the
\texttt{composite\_\allowbreak{}tonelli\_\allowbreak{}shanks} routine.
The algorithm factorises $\sigma(N_L)$ to extract a set of base roots
$\{r_i\}$.  Each root generates an arithmetic progression
$z = r_i + c\,\sigma(N_L)$ for $c \ge 0$, bounded by the global
search limits $z_{\min}$ and $z_{\max}$.

\paragraph{Sigma cache.}
To avoid recomputing $\sigma(p^{2e})$ inside the raycast inner loop,
the engine precomputes a hash-map cache (\texttt{SigmaCache}) keyed by
$(p, 2e)$ for all primes up to 250{,}000 and exponents up to~8.  Cache
misses fall back to the verified Lean backend.

\paragraph{Exact valuation sieving.}
Before any expensive factorisation is attempted on candidates~$z$, the
engine subjects them to rigorous exact sieving using precomputed
illegal valuations
(\texttt{generate\_\allowbreak{}illegal\_\allowbreak{}z\_\allowbreak{}valuations}).
This sieve leverages the global limit (e.g., 250{,}000) to precompute prime power pairs $(p^e, p^{e+1})$ for which
$\sigma(p^{2e})$ violates the mod-8 obstruction
(\cref{thm:sieve-sound}).

For each candidate~$z$, the engine executes a $\mathcal{O}(1)$ exact
divisor state check: taking the remainder
$R \equiv z \pmod{p^{e+1}}$, if $R \equiv 0 \pmod{p^e}$ but
$R \neq 0$, then $\nu_p(z) = e$, implying $\nu_p(N_R) = 2e$ (since
$N_R = z^2$) and hence $p^{2e} \parallel N$.  By
\cref{lem:exact-val}, $\sigma(p^{2e})$ must divide $\sigma(N)$, but
computing $\sigma(p^{2e}) \pmod{8}$ yields an immediate parity
contradiction with \cref{thm:mod8}.  The candidate is proven
structurally impossible and discarded. The low-level execution invariants and verification models for the underlying RNS512 Montgomery multiplication and GCD operations powering these checks are extensively documented in \cref{app:verification}.


%% -------------------------------------------------------------------
\subsection{Lock-Free Concurrency \& Telemetry}
\label{subsec:concurrency_telemetry}

The computational pipeline utilizes the Rust Rayon crate for dynamic,
work-stealing concurrency.  To maximize parallel throughput and avoid
the thundering-herd problem, thread synchronization is managed strictly
via lock-free atomic states and read-write locks, minimizing blocking.
This lock-free architecture not only ensures CPU cores remain fully
saturated but also drives real-time telemetry, allowing external
monitoring tools to track the active primes and search progress
dynamically without introducing observational overhead.

%% -------------------------------------------------------------------
\subsection{Dynamic Minimum Factors}
\label{subsec:dynamic_minimum_factors}

Leveraging our formalized Prasad--Sunitha bound
(\cref{thm:prasad-sunitha}), the computational search floor dynamically
shifts during deep exploration.  When the depth-first search branch
definitively excludes 3 and~5 as prime factors of a candidate~$N$, Lean
calculates the Prasad--Sunitha prime factors (mathematical bounds) natively
within the proof environment. These verified bounds are passed via the FFI
boundary to the Rust execution engine. The Rust engine then uses these bounds
to safely drive its high-performance pruning engine, automatically elevating
the minimum required distinct prime factors from the baseline of~\TelemetryHagisBaselineMinPrimeFactors\
up to~\TelemetryPrasadSunithaBound. This lets the engine reject such branches earlier.


This component heavily relies on our core software dependencies, specifically Lean 4~\cite{moura2021}, Mathlib~\cite{mathlib2020}, Rayon~\cite{rayon2024}, bnum~\cite{bnum2024}, and ed25519-dalek~\cite{ed25519dalek2024}.

%% =================================================================
%%  5.  Results
%% =================================================================
\section{Results}
\label{sec:results}

This section reports what the prototype search engine did and what
it does not establish.  The Lean~4 library proves individual lemmas
that the engine's pruning rules are meant to rely on; for example,
\texttt{rust\_\allowbreak{}sieve\_\allowbreak{}soundness} justifies the
mod-8 rule used by the Rust
\texttt{generate\_\allowbreak{}illegal\_\allowbreak{}z\_\allowbreak{}valuations}
sieve.  No theorem states that the search as a whole is complete, and
the search does not prove any lower bound on quasiperfect numbers.  The
strongest bound we are aware of is $N > 10^{45}$, due to
Alekseyev~\cite{alekseyev2026}; earlier, Hagis and
Cohen~\cite{hagis1982} showed $N > 10^{35}$.


%% -------------------------------------------------------------------
\subsection{Search Coverage and Known Limitations}
\label{subsec:bounds_achieved}

The engine was configured to search up to $10^{\TelemetryMaxLog}$, but
this target is not a result.  The coverage of the search is incomplete
for the following reasons.
\begin{enumerate}
  \item \textbf{Bounded component list.}  Phase~1 only generates prime
    powers $p^{2e}$ with $p < 250{,}000$ and $2e \le 8$
    (\cref{subsec:global_sieve}).  Any QPN with a prime factor of at
    least $250{,}000$, or with a prime appearing to an exponent above~8,
    lies outside the search.
  \item \textbf{Misapplied forced-inclusion rule.}  The cyclotomic
    dependency graph (CDG) ``forced cascade'' treats certain primes as
    components that must appear in~$N$, and prunes a branch when those
    components cannot be added or would push the abundancy past the
    target.  The Lean theorems
    it cites (\texttt{forced\_\allowbreak{}inclusion} and
    \texttt{transitive\_\allowbreak{}forced\_\allowbreak{}inclusion} in
    \texttt{CyclotomicGraph.lean}) show that certain primes divide
    $\sigma(N)$, not that they divide~$N$.  Since
    $\gcd(N, \sigma(N)) = 1$ for a QPN, such primes in fact cannot
    divide~$N$, so this rule has no formal justification and can
    discard valid branches.
  \item \textbf{Pruning bounds relative to the list.}  Several pruning
    bounds, including the suffix-starvation bound, compute the maximum
    achievable abundancy only from components in the engine's finite
    list (\cref{subsec:dfs_prefix}).  They are therefore not valid for
    an arbitrary~$N$ whose factors lie outside that list.
  \item \textbf{Ray-casting did not run.}  The exact ray-casting stage
    (\cref{subsec:ray_casting}) was not exercised; every branch
    recorded in the telemetry was terminated by an earlier pruning
    rule.
\end{enumerate}

For branches whose prefix is disjoint from $\{3, 5\}$ the engine
applies the Prasad--Sunitha bound~\cite{prasad_sunitha}
($\omega(N) \ge \TelemetryPrasadSunithaBound$), which is proved in
Lean.  For prefixes containing both 3 and~5, the bound
$\omega(N) \ge \TelemetryHagisBaselineMinPrimeFactors$ of Hagis and
Cohen is used.

The computation was executed locally.  The hardware environment and physical
execution telemetry metrics are detailed in \cref{tab:hardware_specs}.

\begin{table}[htbp]
    \centering
    \makebox[\textwidth][c]{ \begin{tabular}{|l|l|}
        \hline
        \textbf{Parameter} & \textbf{Specification / Value} \\
        \hline
        Computer / Arch & Apple Mac mini (M4, 10 Cores) \\
        Available RAM & 16 GB Unified Memory \\
        Total Wall-Clock Time & \TelemetryTotalTime \\
        Phase 1 (Precomputation/ECM Factorization) Time & \TelemetryPhaseOneTime \\
        Phase 2 (DFS Traversal) Time & $\approx \TelemetryPhaseTwoTime$ seconds \\
        Phase 1 Eliminations & \TelemetryPhaseOnePruned\ components pruned \\
        Phase 2 Checks & \TelemetryPhaseTwoBranches\ branches evaluated \\
        Run Certificate Hash & \TelemetryCertHash \\
        \hline
    \end{tabular} }
    \caption{Hardware specifications and search telemetry.}
    \label{tab:hardware_specs}
\end{table}

As shown in \cref{tab:hardware_specs}, the main cost
is the up-front Phase~1 integer factorization---specifically computing exact
cyclotomic sums and running the Elliptic Curve Method (ECM) factorization on
tens of thousands of massive $\sigma(p^{2e})$ values before the DFS even starts.
The combinatorial DFS tree traversal (Phase~2) is not the bottleneck; at approximately
\TelemetryNodesPerSec\ nodes per second, it efficiently leverages these precomputations
to process the search space.  Most Phase~2 branches were terminated
by starvation pruning, which, as noted above, is only valid relative to
the engine's component list.  The distribution of branch
terminations is outlined in \cref{tab:pruning_stats}.

\begin{table}[htbp]
  \centering
  \makebox[\columnwidth][c]{ \begin{tabular}{|l|r|}
    \hline
    \textbf{Pruning Mechanism}
      & \textbf{\% of Branches Eliminated} \\
    \hline
    Abundance / Starvation Pruning   & \TelemetryAbundancePct\% \\
    Ray-Casting (Exact Valuations)   &  \TelemetryRaycastPct\% \\
    \hline
  \end{tabular} }
  \caption{Distribution of topological branch pruning strategies across
           DFS execution (Total branches pruned: \TelemetryPruned).}
  \label{tab:pruning_stats}
\end{table}

The telemetry records no runtime panics or overflows during the run.
This shows that the prototype executed without crashing; it does not
show that the search covered every candidate below
$10^{\TelemetryMaxLog}$, for the reasons listed in
\cref{subsec:bounds_achieved}.

%% =================================================================
%%  6.  Conclusion and Future Work
%% =================================================================
\section{Conclusion and Future Work}
\label{sec:conclusion}

The Unified Algebraic-Lattice Bipartition Framework (UALBF) combines a
Lean~4 library of results about quasiperfect numbers with a prototype
Rust search engine.  The Lean library verifies that a QPN is an odd
perfect square~\cite{cattaneo1951}, the modulo-8 obstruction on primes
dividing $\sigma(N)$, the Prasad--Sunitha bound~\cite{prasad_sunitha},
the 2.0442 totient-ratio bound for QPNs coprime to~15, the cyclotomic
and Zsigmondy lemmas, and the multiplicativity of~$\sigma$ over a
prefix--suffix bipartition.  Each of these is a standalone lemma that
can justify a particular pruning step.

The search engine is not a proof.  It only considers primes below
$250{,}000$ with exponents $2e \le 8$, its CDG forced-cascade rule
misreads a Lean theorem about primes dividing $\sigma(N)$ as forcing
primes into~$N$, several of its pruning bounds only account for
components in its finite list, and its exact ray-casting stage did not
run (\cref{subsec:bounds_achieved}).  No theorem states that the search
as a whole is complete, and we claim no lower bound on quasiperfect
numbers; the strongest bound we are aware of is
$N > 10^{45}$~\cite{alekseyev2026}.
As detailed in \cref{app:tcb} and \cref{app:verification}, the
Lean-to-Rust FFI boundary is also unverified and forms part of the
Trusted Computing Base.

Future work will focus on the following:

\begin{enumerate}
  \item \textbf{Soundness of the Search:}
    Removing or correctly justifying the CDG forced-cascade rule,
    making every pruning bound valid for components outside the
    engine's list, handling primes above the sieve limit and larger
    exponents, and stating in Lean what a completed search would
    prove.

  \item \textbf{Deep Formalisation of Divisibility Chains:}
    Formalising generalised modulo-3 and modulo-9 constraints in
    Lean~4 to further restrict permissible suffix structures for
    candidates divisible by~3, where the Prasad--Sunitha bound does
    not apply.

  \item \textbf{Distributed Scaling:}
    Expanding the Rust engine's architecture to a distributed cluster,
    allowing distinct algebraic ``cubes'' to be assigned dynamically
    across hundreds of nodes.

  \item \textbf{Enhanced Lattice Pruning:}
    Augmenting the sieving with tighter tolerance analysis to prune starvation branches at
    shallower depths.
\end{enumerate}

Only after these gaps are closed could a run of the engine support a
claim about the non-existence of quasiperfect numbers in a range.

\appendix
\section{Trusted Computing Base (TCB) \& Verification Boundaries}
\label{app:tcb}

To transparently align the formal claims of this manuscript with the software's actual state, we explicitly define the Trusted Computing Base (TCB) boundaries below. The following components are treated as unverified mathematical assumptions or trust boundaries within the codebase, rather than completed proofs:

\begin{enumerate}
    \item \textbf{Unverified FFI Boundary:} The Foreign Function Interface (FFI) boundary between the Lean 4 formalization and the Rust execution engine is unverified. While the Rust engine trusts the semantics exported via Lean's \texttt{@[export]} pragmas (such as the C-compatible data serialization), the bridging logic itself is not mechanically checked and forms a boundary in the TCB.
    
    \item \textbf{Inactive GPU Pipeline Status:} The highly parallelized batch-factorization GPU Pollard's Rho pipeline, implemented in Apple Metal, is completely bypassed in the active verified paths. High-performance execution relies entirely on sequential CPU loops, leaving the GPU pipelines unverified and unused in QPN search invariants.

    \item \textbf{Search Completeness Not Proven:} The Lean library proves individual pruning lemmas, but no theorem states that the search as a whole covers every candidate in its target range. The known coverage gaps are listed in \cref{subsec:bounds_achieved}.
\end{enumerate}

\section{Exhaustive Verification of Low-Level Computational Primitives}
\label{app:verification}

This appendix provides exhaustive technical evidence of the formal verification methodologies applied to the low-level execution invariants of the computational engine. In alignment with the stratified architecture defined in Section~2, our goal is to bridge abstract Lean~4 proofs with high-performance execution without compromising mathematical rigor.

\subsection{Pocklington Primality Testing Verification}
\label{app:subsec:miller_rabin}

The computational engine heavily relies on deterministic primality testing for evaluating component feasibility, specifically via the \texttt{verified\_is\_prime} routine. This function performs a fully verified Pocklington primality test, completely eliminating unverified mathematical assumptions.

To specify the theoretical framework, the engine mathematically links the Fermat condition and the GCD condition via \texttt{lemma\_pocklington\_certificate}. By establishing absolute mathematical certainty through automated, mechanical proofs in Verus, the overarching mathematical engine runs a zero-assumption primality pipeline. This formally proved pipeline reduces the TCB footprint for numbers under $2^{64}$ and beyond.

\subsection{RNS512 Arithmetic Models and GPU Execution}
\label{app:subsec:rns512}

For deeper search horizons requiring 512-bit arithmetic (such as Pollard's rho factorization and raycast sieving), the engine employs a Residue Number System (RNS) architecture, specifically \texttt{Rns512}, represented as an 8-channel array of 64-bit integers.

Safely executing 512-bit arithmetic on Highly Parallel GPU architectures introduces significant complexity. To secure this layer, we specify formal models for the core RNS512 operations---specifically Montgomery multiplication and Greatest Common Divisor (GCD) algorithms. 

The Montgomery multiplication specification mathematically models the invariants of the multi-precision reduction sequence, guaranteeing that intermediate modular additions and subtractions preserve congruence without silent overflow across the 64-bit boundaries. Similarly, the GCD specification formalizes the termination and correctness of the Euclidean steps applied to the \texttt{Rns512} channels.

\subsection{Verification Methodology}
\label{app:subsec:hybrid_verification}

Given the challenges of deploying formally verified code directly to GPU environments (e.g., Metal kernels), and because Pollard's Rho GPU acceleration is currently categorized as \textbf{Inactive} (and restricted strictly to macOS/Metal environments), the framework adopts a CPU-bound testing methodology to ensure the safety and correctness of the underlying computational logic.

This methodology stratifies the verification into two interdependent layers:
\begin{enumerate}
    \item \textbf{Verus Formal Proofs:} We employ Verus to formally verify the structural logic and algebraic properties of the CPU-bound reference implementations (such as the RNS512 models and the computational structure of the deterministic Pocklington primality tests, removing any reliance on unverified sufficiency assumptions). This provides a certified bedrock for the execution logic.
    \item \textbf{Property-Based Testing:} To validate the conceptual models of GPU operations, we employ property-based testing (via \texttt{proptest}) entirely on the CPU side. By subjecting the placeholder Montgomery multiplication and GCD logic to 10,000 generated test cases asserting structural parity properties (e.g., \texttt{test\_mont\_mul\_parity\_property} and \texttt{test\_gcd\_parity\_property}), we establish foundational confidence in the logic intended for future GPU execution.
\end{enumerate}

This strategy tightly couples rigorous formal verification with practical fuzzing, delivering exhaustive mathematical confidence and empirical runtime correctness for the active CPU search engine.

\subsection{Formal Verification Manifest}
\label{app:subsec:verification_manifest}
This section dynamically lists the extracted semantic hashes from Lean theorems and Rust implementation functions, confirming the cryptographically verifiable linkage between the mathematical formalisms and high-performance execution.

\input{verification_manifest.tex}


\bibliographystyle{alpha}
\bibliography{references}


\end{document}
